% !TeX root = ../papers/nc-p3.tex
\section{二次曲面間の双有理対応と孤立ファイバーを許す曲面被覆}
\label{sec:birational-quadric-chains}
\begin{theorem}\label{thm:birational-quadric-chains}
標準正則窓の点対応に、全整数位置で整な二次曲面から次の整な二次曲面への支配的双有理グラフが含まれるとする。
するとこれらの双有理写像は全て、周囲の射影三空間の同型の制限である。
余分な点対や正次元の特殊ファイバーを初めから除く必要はない。
従ってこの全域の窓は既出の滑らかな二次曲面商または二次錐商のリストに入る。
\end{theorem}

\begin{proof}[一般超平面切断と母線束]
辺を $Q\dashrightarrow Q'$ とする。$Q$ の一般超平面切断は滑らかな二次曲線であり、
基点を避け、像曲線に双有理に写る。
像が平面に入れば次数は2以下である（整二次曲面の平面切断に入る）。
非平面ならこの二次曲線上の点座標を二項二次の完全系 $S^2$、相手を基点を除いた
四空間 $W\subset S^d$ とする。定理\ref{thm:branching-centers-curve}の一般平面の補題により
六つの双一次式は線型行列として独立のままであり、全て乗法 $S^2\otimes W\to S^{d+2}$ の核に入る。
従って $\dim S^2 W\le6$。二回の二項増加不等式はどちらも等号であり、
定理\ref{thm:branching-lines-only}の最小増加の議論から $W=S^3$、$d=3$ である。

有理写像の四切断から固定因子を除いた可動系を取る。
滑らかな $Q$ ではその型 $(a,b)$ は $a,b\ge1$, $a+b\le3$ である。
実際、その一般超平面切断との交叉数が先の像曲線の次数である。
錐では解消 $\mathbb F_2$ の例外切断 $E$ と母線 $F$ を使い、超平面を $E+2F$、
可動系を $aE+bF$ と書く。支配性で $a\ge1$、$E$ が固定因子でないことから $b\ge2a$、
また $b=(E+2F)(aE+bF)\le3$ である。
どちらの場合も射影同型でない写像は、少なくとも一方の母線族を相手の母線族へ射影同型で運ぶ。
実際、$Q'$ が滑らかなら二つの射影直線への可動鉛筆の和がこの可動系である。
これは二つの基点自由な有理座標比の積の四座標を取り、各比の固定因子を除いても
積系には新たな固定因子がないことによる。
適切な源母線への次数の和が1なので、一方の鉛筆はその母線で一定である。
その基底写像の次数は1である。一般点での一対一性と母線から相手母線への非定数性による。
$Q'$ が錐なら、頂点からの投影を通じる二項鉛筆を $P$ と書くと、相手の底三座標の系は
可動部分 $2P$ と有効固定部分の和になる。従って源の同じ母線で $2P$ の次数は高々1、
この鉛筆は母線上で一定である。錐がどちらかにある場合、その側の族は唯一である。
滑らかな両側で両族とも対応すれば各基底直線の同型の積で既に射影同型である。

対応する一つの母線の基底直線を同定する。各位置で母線上の一次切断を基底直線上の束 $E_i$ にまとめる。
滑らかな場合 $E_i=\mathcal O(1)^2$、錐では $E_i=\mathcal O(2)\oplus\mathcal O$ である。
どちらも全切断は四文字 $V_i$、評価 $V_i\otimes\mathcal O\to E_i$ の核 $N_i$ は $\mathcal O(-1)^2$ である。
後者は完全二次系の二つの一次関係からも分かる。
相手母線への写像は基礎体 $k(s)$ 上の射影直線の同型なので、
そのグラフを消す双一次式は非退化な一本である。その直線を
$E_i\otimes E_{i+1}$ で飽和して $J_i\simeq\mathcal O(j)$ と書く。
六関係を二位置で同じ基底 $s$ に沿って評価すれば $H^0(J_i)$ に入る。
全二文字の制限は $H^0(E_i\otimes E_{i+1})$ へ全射で、この空間は次元12である
（分裂の各対で二項形式の乗法が全射）。従って $j\ge1$ である。
一方、$J_i\subset E_i\otimes E_{i+1}$ の行列式は次数 $4-2j$ の非零切断なので $j\le2$。
$j=2$ ならこれは定数非零で、$E_i^*\otimes\mathcal O(2)\simeq E_{i+1}$ を与える。
階数2なのでこれは両 $E$ の同型でもある（行列式を固定すれば $E_i^*(2)\simeq E_i$）。
全切断と各母線上の核直線の対応により、曲面間の写像は射影線型同型に延びる。
$j=1$ なら、四次元の制限核全体が $R_i$ に入り、$R_i$ の像は $H^0(J_i)$ 全部である。
\end{proof}

\begin{proof}[母線を交換した二辺の三次元数]
射影同型でない二辺が一つの滑らかな中間曲面で異なる母線を用いることはない。
そのような二辺を $0,1$ とし、中間四文字を $V_1=B\otimes T$, $\dim B=\dim T=2$ とする。
第一辺は $B$ の基底直線、次の辺は $T$ の基底直線で $j=1$ の表示を持つ。
まず双方の四次元制限核だけを三文字で割ると
\[
          H^0(E_0(1))\otimes H^0(E_2(1))
\]
を得て次元は36である。実際、第一の核は $V_0\otimes B\to H^0(E_0(1))$ の核に
自由な $T$ を掛けたもの、第二も $T,V_2$ について同じで、この二商は別因子で取れる。
$J_0\hookrightarrow E_0\otimes\mathcal O(1)\otimes T$ の係数は
$P\in H^0(E_0)\otimes T$、$J_1$ の係数は $Q\in B\otimes H^0(E_2)$ である。
残る二関係の挿入は、第一辺では $P$ を $B\otimes V_2$ の任意の切断と掛ける像、
第二辺では $Q$ を $V_0\otimes T$ の任意の切断と掛ける像となる。
従って二つの像の次元は各々8以下で、共通に積 $P Q$ を含む。
この積は零でない。一般の二つの基点を選び、それぞれの一次形式因子を評価した
$P,Q$ の値は両方非零であり、残る二つの束因子の普通のテンソル積となるからである。
従って三文字商は $36-8-8+1=21$ 次元以上となり、標準値20に反する。
\end{proof}

\begin{proof}[同じ母線を使う隣接辺]
いま射影同型でない辺 $i$ を一つ取る。両隣が射影同型ならその母線をそのまま移せる。
非射影的な隣辺も、上の三次元数と錐での母線の一意性により、同じ母線を使う必要がある。
従って四つの連続位置 $i-1,i,i+1,i+2$ の基底直線を同定でき、両隣辺の一般母線間の写像は同型である。
標準四テンソルの中央二位置をこの共通母線上で制限すると、$J_i$ の包含と
\[
         H\in H^0\big(\mathcal O(1)\otimes V_{i-1}\otimes V_{i+2}\big)
\]
の積になる。実際 $w$ の中央二位置の縮約像は $R_i$ であり、その制限が $H^0(J_i)$ である。
一般の基底点 $s$ を固定し、右隣の点を同じ母線上で任意に与える。
その前点を位置 $i+1$ に取り、位置 $i$ の母線上では中央の非退化な双一次式を非零にできる。
隣接点対で $w$ は消えるため、$H(s)$ は指定した右隣点を消す。同様に左隣の母線全体を消す。
従って $H$ は $\mathcal O(1)\otimes N_{i-1}\otimes N_{i+2}$ の大域切断である
（一般点での包含は部分束への包含なので全点に延びる）。
ところがこの束は $\mathcal O(-1)^4$ で非零切断がない。
これは $R_i$ の制限が二次元であることに反する。
ゆえに全辺は射影同型であり、独立な基底変更で曲面座標環の反復への全射が得られる。
定理\ref{thm:smooth-quadric-quotients}または定理\ref{thm:quadric-cone-quotients}を適用する。
\end{proof}

\begin{corollary}\label{cor:pure-free-surface-complete}
どの関係空間にも純テンソルがない標準正則窓について、任意の辺の射影像が曲面以上なら、
全窓は可換多項式窓、正則正規一次直線の平面商、滑らかな二次曲面商、または整二次錐商の
既出の添字保持分類で尽きる。点対応の全射影が点ファイバーという仮定は不要である。
\end{corollary}
\begin{proof}
三次元像は定理\ref{thm:branching-lines-only}で済む。
純関係がないので全辺で直線ファイバーの中心は有限である。
面上の一般点は隣接辺で曲線に縮約できない。縮約する曲線の一般点には直線ファイバーが生じるからである。
従って与えた面の一般点の相手は再び面の一般点であり、同じ辺を戻れば一意である。
この二点で隣辺の評価階数は一般に3以上で、隣接切断の直交不等式によりともに3となる。
よってこの対応は前後に延び、面の間の双有理写像の列を得る。
面は既に平面または整二次曲面に限られる。

平面が一つあれば、その次も前も平面で写像は射影同型である。
実際、平面への辺の六関係の制限は純関係の不在で独立である。
双有理写像の像が全三空間を張れば、相手の四座標に源の三一次式を掛けた包は次元6以下になる。
しかし四つの異なる先頭単項式の三文字倍の異なるものは7個以上である。
$12-6=6$ 個しかなければ全対が異なる一文字倍で会うので、四単項式は同一の単項式と
四つの異なる文字の積となる。三文字しかなく不可能である。
相手も平面なら、さらに相手を三文字へ射影して失う関係は三次元以下で、独立な線型関係が三つ以上残る。
共通因子を除いた三座標を $f$ とする。独立な二つの一次関係は有理関数体上も独立である。
従属なら共通の定数関係の一次式倍となるが、像は平面を張る。二関係の外積は $f$ の多項式倍なので
$\deg f\le2$。二次なら各外積は $f$ の定数倍で、クラメル則により第三の関係も最初の二つの定数結合となり矛盾する。
従って写像は一次である。
従って平面対上の制限関係はちょうど三つの可換関係となる。

平面方程式 $a_i$ を取ると、この射影の三次元核は
$a_i\otimes V_{i+1}+V_i\otimes a_{i+1}$ に入る。
両混合三空間 $a_i\otimes(V_{i+1}/ka_{i+1})$ と $(V_i/ka_i)\otimes a_{i+1}$ のいずれへも同型である。
実際、一方への射影で零なら元の関係は純となる。
従って $a_i V_{i+1}=V_i a_{i+1}$、全位置で正規となり、商は三変数多項式窓である。
元と商のHilbert関数から両側の乗法は単射であり、定理\ref{thm:plane-quotient-unmarked}の分類に入る。
平面がなければ全位置で整二次曲面となり、前定理を適用できる。
\end{proof}
